% Standalone resolution manuscript generated from solution.md.
% Compile with XeLaTeX; author and review metadata are maintained in Markdown.
\documentclass[10pt,a4paper]{article}
\usepackage[margin=24mm,headheight=15pt]{geometry}
\usepackage{fontspec}
\setmainfont{DejaVuSerif.ttf}[BoldFont=DejaVuSerif-Bold.ttf,ItalicFont=DejaVuSerif-Italic.ttf,BoldItalicFont=DejaVuSerif-BoldItalic.ttf]
\setsansfont{DejaVuSans.ttf}[BoldFont=DejaVuSans-Bold.ttf,ItalicFont=DejaVuSans-Oblique.ttf]
\setmonofont{DejaVuSansMono.ttf}
\usepackage{amsmath,amssymb,unicode-math}
\setmathfont{latinmodern-math.otf}
\usepackage{xcolor,microtype,parskip,needspace,fancyhdr}
\usepackage{bookmark}
\usepackage{xurl}
\hypersetup{colorlinks=true,urlcolor=blue,linkcolor=blue,pdftitle={Parity
obstructions for the conditioning of sign matrices},pdfauthor={Matthew
J. Colbrook},pdflang={en-GB}}
\urlstyle{same}
\setlength{\emergencystretch}{3em}
\providecommand{\tightlist}{\setlength{\itemsep}{0pt}\setlength{\parskip}{0pt}}
\providecommand{\pandocbounded}[1]{#1}
\setcounter{secnumdepth}{0}
\allowdisplaybreaks[1]
\widowpenalty=10000
\clubpenalty=10000
\pagestyle{fancy}
\fancyhf{}
\fancyhead[L]{\small\sffamily PARTIAL RESULT / IS-05}
\fancyhead[R]{\small\sffamily 11 September 2026}
\fancyfoot[L]{\parbox[b]{0.9\textwidth}{\fontsize{8}{10}\selectfont Independent
Codex-agent verification; no external human peer review or formal
certification.}}
\fancyfoot[R]{\footnotesize\thepage}
\begin{document}
{\raggedright\Large\bfseries Parity obstructions for the conditioning of
sign matrices\par}
\vspace{0.4em}
{\normalsize\bfseries Matthew J. Colbrook\par}
{\small Department of Applied Mathematics and Theoretical Physics,
University of Cambridge, Cambridge, United Kingdom\par}
{\small\href{mailto:m.colbrook@damtp.cam.ac.uk}{m.colbrook@damtp.cam.ac.uk}\par}
\vspace{0.6em}
\textbf{Status of this manuscript:} Independently checked partial
result; the exact exponent remains open.\\
\textbf{Target:}
\href{https://github.com/ajt60gaibb/OpenProblemsInNLA/blob/b4123194697bdf6f8f82518c1dd7d6c40a30c2e0/eigenvalues-and-inverse-problems/IS-05/README.md}{IS-05,
original catalog snapshot}.\\
\textbf{Reviewed:} 11 September 2026.

The original proof draft was generated in a ChatGPT conversation. A
separate Codex agent independently checked the complete argument against
the exact target; its
\href{https://github.com/ajt60gaibb/OpenProblemsInNLA/blob/main/references/colbrook-additional-2026-09-11/verification/reviews/IS-05-review.md}{detailed
review} records \textbf{PASS} for the partial bound only. The
mathematical sections below are unchanged from the reviewed draft. This
is independent agent verification, not external human peer review or a
formal proof certificate. The
\href{https://github.com/ajt60gaibb/OpenProblemsInNLA/blob/main/references/colbrook-additional-2026-09-11/README.md}{submission
record} preserves the original manuscript and verification history.

\section{Target and result}\label{target-and-result}

The repository defines \[
 h(n)=\min_{A\in\{-1,1\}^{n\times n}}\kappa_2(A),
\] where singular matrices have infinite condition number, and asks for
the exact value \[
 \alpha_*=\sup\{\alpha\ge0:\exists C>0\ \forall n\ge1,
            \ h(n)-1\le Cn^{-\alpha}\}.
 \tag{1}
\] Its checked formulation records \(17/92\le\alpha_*\le1\).\footnote{\emph{Open
  Problems in Numerical Linear Algebra}, IS-05, ``The optimal decay
  exponent for the conditioning of sign matrices,'' checked 11 September
  2026;
  \href{https://github.com/ajt60gaibb/OpenProblemsInNLA/blob/main/eigenvalues-and-inverse-problems/IS-05/README.md}{canonical
  entry}.} This note improves the upper bound in that displayed
interval, without resolving the exact-value question.

\textbf{Theorem 1 (odd-order obstruction).} If \(n\) is odd, every
nonsingular sign matrix \(A\in\{-1,1\}^{n\times n}\) satisfies \[
 \kappa_2(A)\ge
 \sqrt{1+\frac{n-1}{n^2}}+\frac{\sqrt{n-1}}{n}.
 \tag{2}
\] Consequently, \[
 h(n)-1\ge\frac{\sqrt{n-1}}n\quad(n\text{ odd}),
 \qquad \boxed{\alpha_*\le\frac12}.
 \tag{3}
\]

\section{A variance bound in terms of the condition
number}\label{a-variance-bound-in-terms-of-the-condition-number}

Let \(G=A^TA\), let \(m,M\) be its smallest and largest eigenvalues, and
put \(t=M/m=\kappa_2(A)^2\). Since every diagonal entry of \(G\) equals
\(n\), its eigenvalues \(\lambda_1,\ldots,\lambda_n\) have mean \(n\).
Define their variance \[
 v=\frac1n\sum_{i=1}^n(\lambda_i-n)^2
   =\frac1n\operatorname{tr}(G-nI)^2.
 \tag{4}
\] Summing the nonnegative quantities \((M-\lambda_i)(\lambda_i-m)\)
yields \[
 v\le(M-n)(n-m)=(tm-n)(n-m).
\] As a quadratic function of \(m\), the last expression is at most its
unrestricted maximum, \[
 v\le\frac{n^2(t-1)^2}{4t}.
 \tag{5}
\] Writing \(\kappa=\kappa_2(A)\ge1\) and taking square roots gives \[
 \kappa-\kappa^{-1}\ge\frac{2\sqrt v}{n},
 \qquad
 \kappa\ge\sqrt{1+\frac v{n^2}}+\frac{\sqrt v}{n}.
 \tag{6}
\] This derivation also covers \(t=1\); then (5) forces \(v=0\).

\section{Parity supplies the necessary
variance}\label{parity-supplies-the-necessary-variance}

For odd \(n\), the inner product of two sign vectors of length \(n\) is
an odd integer, so every off-diagonal entry \(g_{ij}\) of \(G\) has
\(g_{ij}^2\ge1\). Therefore \[
 v=\frac1n\sum_{i\ne j}g_{ij}^2\ge n-1.
 \tag{7}
\] Substituting (7) into (6) proves (2), and hence the first part of
(3).

Suppose some \(\alpha>1/2\) were admissible in (1). For every odd \(n\),
\[
 C\ge n^\alpha(h(n)-1)
   \ge n^{\alpha-1}\sqrt{n-1}.
\] The right-hand side diverges along the odd integers, a contradiction.
This proves the upper bound on \(\alpha_*\). Singular matrices cause no
exception: their infinite condition numbers satisfy the lower bounds
automatically and do not improve the minimum.

\section{An additional bound for orders equal to two modulo
four}\label{an-additional-bound-for-orders-equal-to-two-modulo-four}

\textbf{Theorem 2.} If \(n\equiv2\pmod4\), every nonsingular sign matrix
of order \(n\) satisfies \[
 \kappa_2(A)\ge
 \sqrt{1+\frac{2(n-2)}{n^2}}+\frac{\sqrt{2(n-2)}}n.
 \tag{8}
\]

\textbf{Proof.} Form a graph whose vertices are the columns of \(A\),
with an edge for each pair of orthogonal columns. This graph is
triangle-free. Indeed, if three sign columns were pairwise orthogonal,
multiply their rows by signs to make the first column all ones. The
second has \(n/2\) plus entries and \(n/2\) minus entries. Orthogonality
of the third to the first two forces its sum on each of these two halves
to be zero. Each half must consequently have even size, implying
\(4\mid n\), a contradiction.

A triangle-free graph on \(n\) vertices has at most \(n^2/4\) edges. To
see this directly, for every edge \(uv\) its endpoint degrees satisfy
\(d_u+d_v\le n\). Summing over all \(e\) edges gives
\(\sum_vd_v^2\le ne\), whereas Cauchy--Schwarz gives
\(\sum_vd_v^2\ge(2e)^2/n\). Thus \(e\le n^2/4\) (including \(e=0\)).

At least \[
 \frac{n(n-1)}2-\frac{n^2}4=\frac{n(n-2)}4
\] unordered column pairs therefore have nonzero inner products. These
inner products are even integers and have square at least four. Counting
both orders in (4) gives \[
 v\ge\frac2n\,4\,\frac{n(n-2)}4=2(n-2).
\] Equation (6) now proves (8).

\section{What remains open}\label{what-remains-open}

Combining Theorem 1 with the lower bound quoted by the repository gives
the proposed updated interval \[
 \frac{17}{92}\le\alpha_*\le\frac12.
\] No construction establishing the matching exponent \(1/2\) is
supplied here. Theorems 1 and 2 are therefore partial progress, not a
solution to the exact-value target. The repository submission records
\textbf{Partially resolved}, with the exact exponent still open.

The verification script checks the finite-dimensional inequalities on
small sign matrices. The all-orders proof is (4)--(8), not the finite
tests. Whether these elementary bounds have appeared elsewhere remains
to be checked before any novelty claim.

\section{Sources}\label{sources}

\begin{enumerate}
\def\labelenumi{\arabic{enumi}.}
\item
  \emph{Open Problems in Numerical Linear Algebra}, IS-05,
  \href{https://github.com/ajt60gaibb/OpenProblemsInNLA/blob/main/eigenvalues-and-inverse-problems/IS-05/README.md}{canonical
  statement}, checked 11 September 2026.
\item
  B. Alexeev, J. Jasper, and D. G. Mixon, \emph{Asymptotically optimal
  approximate Hadamard matrices},
  \href{https://arxiv.org/html/2511.14653v1}{arXiv:2511.14653v1}, §6,
  Problem 11; the primary reference identified by the repository for the
  exponent question and its quoted bounds.
\item
  The later \href{https://arxiv.org/html/2511.14653v2}{arXiv version 2},
  dated 18 August 2026, retains the construction exponent below
  \(17/92\) and the earlier obstruction in §5, Problem 11. The
  independent review checked this current source; the new upper-bound
  proof above is unchanged.
\end{enumerate}
\end{document}
